\begin{python0}
from solutions import *; clear()
\end{python0}
\sectionthree{Why inline functions?}

Why do we have inline functions?

Like I said it's faster (see caveat later). Why?

Because there is some CPU cost in just making a function call. Look at the notes on Functions. Recall that to make a function call, the CPU must set up a stack which contains at least information on how to go back to the caller function. That takes time. It's the same when a function call returns. In other words, just \EMPHASIZE{the act of calling the function and returning from a function call} (i.e. not including the computation done in the body of function called) \textbf{takes time}. In fact it's even possible that the cost of making a function call and returning from the call takes more time than the actual computation in the body of the function!!!

Inline functions calls do \textbf{\textit{not}} involve function calls at all because the inline function's body has already been copied to the place where you make the function call.


